% Options for packages loaded elsewhere
\PassOptionsToPackage{unicode}{hyperref}
\PassOptionsToPackage{hyphens}{url}
\documentclass[
  a4paper,
]{article}

\usepackage{xcolor}
\usepackage{amsmath,amssymb}

\usepackage{ctex}
\usepackage{xeCJK}
\setCJKmainfont{Noto Serif CJK SC}[BoldFont=Noto Sans CJK SC,ItalicFont=LXGW WenKai]

\DeclareMathOperator{\tr}{tr}

\setcounter{secnumdepth}{-\maxdimen} % remove section numbering

\usepackage{bookmark}
\IfFileExists{xurl.sty}{\usepackage{xurl}}{} % add URL line breaks if available
\urlstyle{same}

\usepackage{iftex}
\defaultfontfeatures{Scale=MatchLowercase}
\defaultfontfeatures[\rmfamily]{Ligatures=TeX,Scale=1}
\usepackage{lmodern}

\makeatletter
\@ifundefined{KOMAClassName}{% if non-KOMA class
  \IfFileExists{parskip.sty}{%
    \usepackage{parskip}
  }{% else
    \setlength{\parindent}{0pt}
    \setlength{\parskip}{6pt plus 2pt minus 1pt}}
}{% if KOMA class
  \KOMAoptions{parskip=half}}
\makeatother

\setlength{\emergencystretch}{3em} % prevent overfull lines
\providecommand{\tightlist}{\setlength{\itemsep}{0pt}\setlength{\parskip}{0pt}}

\usepackage{bookmark}
\IfFileExists{xurl.sty}{\usepackage{xurl}}{} % add URL line breaks if available
\urlstyle{same}
\hypersetup{hidelinks,pdfcreator={cjx}}


\usepackage{titling}
\title{期中考试}
\author{}
\date{2026年5月14日}

\usepackage{geometry}
\geometry{a4paper,left=1.27cm,right=1.27cm,top=1.7cm,bottom=1.7cm}

\usepackage{enumitem}
\usepackage{etoolbox}
\AtBeginEnvironment{enumerate}{
  \apptocmd{\tightlist}{\def\labelenumii{(\arabic{enumii})}}{}{}
}


\usepackage{fancyhdr}

\begin{document}
\maketitle

\pagestyle{fancy}
\fancyhf{}

\lhead{2026年5月14日}
\chead{期中考试}

\renewcommand{\headrulewidth}{1pt}
\rfoot{\thepage}

\begin{enumerate}
\def\labelenumi{\arabic{enumi}.}
\tightlist
\item
  判断\(\mathbb{R}\)上的矩阵\(A\)和\(B\)是否可以对角化. \[
   A = \begin{pmatrix}
   1 & 6 & & \\
   & 1 & 7 & \\
   & & 5 & \\
   & & & 5
   \end{pmatrix},
   B = \begin{pmatrix}
   1 & 6 & & \\
   & 1 & 7 & \\
   & & 5 & 7 \\
   & & & 5
   \end{pmatrix}
   \]
\item
  计算以下矩阵的行列式. \[
   \begin{pmatrix}
   1 & 1 & 1 & 1 \\
   2 & 1 & 1 & -3 \\
   1 & 2 & 2 & 5 \\
   4 & 3 & 2 & 1 
   \end{pmatrix}
   \]
\item
  \begin{enumerate}
  \def\labelenumii{\arabic{enumii}.}
  \tightlist
  \item
    陈述方阵\(A\)的特征多项式(以下均指首一特征多项式)\(\chi_A\)的定义.
  \item
    陈述方阵\(A\)的极小多项式(以下均指首一极小多项式)\(\pi_A\)的定义.
  \item
    利用上述两个多项式陈述 Cayley-Hamilton 定理.
  \item
    令\(A\)为\(\mathbb{Q}\)系数方阵.
    假设\(A\)的特征多项式\(\chi_A = T^7 - 2T^6 + 3T^5 - 4T^4 + 5T^3 - 6T^2 + 7T - 8\).
    求\(\det A\)和\(A^{-1}\)的特征多项式\(\chi_{A^{-1}}\).
  \item
    令\(A\)为方阵.
    求证极小多项式\(\pi_A\)的常数项\(\pi_A(0)\)非零当且仅当\(A\)可逆.
  \item
    令\(A\)为可逆方阵,
    且极小多项式为\(\pi_A = T^5 + 137T^3 - 47T^2 + 2\),
    求\(A^{-1}\)的极小多项式\(\pi_{A^{-1}}\)(只需给出答案不要求证明).
  \item
    令\(A\)为可逆方阵.
    从上题中总结出\(A\)的极小多项式\(\pi_A\)与\(A^{-1}\)的极小多项式\(\pi_{A^{-1}}\)之间的一个关系式,
    并给出证明.
  \end{enumerate}
\item
  令\(K\)为一个域且包含于\(\mathbb{C}\)之中.

  \begin{enumerate}
  \def\labelenumii{\arabic{enumii}.}
  \tightlist
  \item
    令\(A \in \mathcal{M}_n(K)\).
    说明\(\exp A = \sum\limits_{k = 0}^{+\infty} \frac{1}{k!} A^k \in \mathcal{M}_n(\mathbb{C})\)是良好定义的.
    即说明为什么这个无穷级数是收敛到一个矩阵(但要注意矩阵的元素不一定仍然在\(K\)中).
  \item
    令\(A \in \mathcal{M}_n(K)\). 求证: \(\det \exp A = \exp \tr A\).
  \item
    令\(A \in \mathcal{M}_n(\mathbb{R})\)以及\(M = \exp A\), 求证:
    \(M \in \mathcal{M}_n(\mathbb{R})\)且\(\det M > 0\).
    并举例说明存在\(M \in \mathcal{M}_2(\mathbb{R})\) (即\(n = 2\))
    使得\(\det M > 0\)但是对任意\(A \in \mathcal{M}_2(\mathbb{R})\)均有\(M \neq \exp A\).
    提示: \(\exp A\)总是可以''开平方''的.
  \item
    令\(M \in \mathcal{M}_n(\mathbb{C})\)为幂零矩阵,
    是否存在\(A \in \mathcal{M}_n(\mathbb{C})\)使得\(M = \exp A\).
    为什么? 若存在则给出例子, 若不存在则给出证明.
  \item
    求出矩阵\(A = \begin{pmatrix} 4 & -14 & 8 \\ 1 & -5 & 4 \\ 2 & -5 & 1 \end{pmatrix}\)的
    Jordan 标准型, 并给出过渡矩阵以及其逆.
  \item
    上述矩阵\(A\)的极小多项式是什么?
  \item
    \(A\)同上题, 给出\(A^n\)关于\(n\)的公式.
  \item
    \(A\)同上题, 求出\(\exp A\).
  \end{enumerate}
\end{enumerate}

\end{document}
